% GENERATED FILE. Edit canonical lessons, then run npm run generate:books.
% canonical-source-hash: fnv1a64:e64649467d8b3fe6
% canonical-lessons: ES-C477-ahorro, ES-C477-alquiler, ES-C477-multa, ES-C477-ganar, ES-C477-trescientos

\chapter{El dinero de cada mes}
\label{ch:es-el-dinero-de-cada-mes}
\hlchaptermodality{477}

\begin{chapteropening}
\textbf{By the end of this chapter:} \emph{I can read a note about savings, rent, a fine and what someone earns.}

Complete the last lesson of chapter 477: I can read a note about savings, rent, a fine and what someone earns.
\end{chapteropening}

\section[el ahorro]{el ahorro — the saving; savings}

\label{lesson:ES-C477-ahorro}



\begin{quote}
Before the new one: say the Spanish for the mistake, then the Spanish for bored.
\end{quote}

\subsection*{You'll want to know: el ahorro}
\textbf{el ahorro} — \textquotedblleft{}the saving; savings\textquotedblright{}.

\begin{quote}
Tengo unos ahorros para el verano.
\end{quote}

\begin{quote}
\emph{I have some savings for the summer.}
\end{quote}

In the plural, \emph{los ahorros} is the money you have put aside.

\begin{cousinweb}
\textbf{ahorro} is the noun of \emph{ahorrar}, \textquotedblleft{}to save\textquotedblright{}. \emph{Ahorrar} is built on \emph{horro}, \textquotedblleft{}freed\textquotedblright{}, said of a slave set free, from Arabic \emph{ḥurr}, \textquotedblleft{}free\textquotedblright{} (EtymDB 2.1): to save was first to free, then to spare.
\end{cousinweb}

\begin{grammarlens}[title={What you've built}]
One more word for reading notices, forms and messages.
\end{grammarlens}

\subsection*{Your turn}
\begin{itemize}
\raggedright
  \item \emph{Say it:} \emph{el ahorro}
  \item \emph{Say it:} \emph{el ahorro}, once more
  \item \emph{Say it:} a sentence of your own with \emph{el ahorro}
  \item \emph{Recall:} say the Spanish for the mistake, then the Spanish for bored, then say \emph{el ahorro} again
\end{itemize}

\begin{tcolorbox}[breakable,colback=teal!4,colframe=teal!35!black,title={Before you move on}]
What does el ahorro mean? (\textquotedblleft{}The saving; savings\textquotedblright{}.) Say it once more.
\end{tcolorbox}
\section[el alquiler]{el alquiler — the rent; renting}

\label{lesson:ES-C477-alquiler}



\begin{quote}
Before the new one: say the Spanish for bored, then the Spanish for the saving.
\end{quote}

\subsection*{You'll want to know: el alquiler}
\textbf{el alquiler} — \textquotedblleft{}the rent; renting\textquotedblright{}.

\begin{quote}
El alquiler sube en enero.
\end{quote}

\begin{quote}
\emph{The rent goes up in January.}
\end{quote}

\emph{Se alquila} on a sign means for rent.

\begin{cousinweb}
\textbf{alquiler} $\leftarrow$ Hispanic Arabic \emph{alkirí}, from Arabic \emph{al-kirāʾ}, \textquotedblleft{}the hire, the rent\textquotedblright{} — the Arabic article \emph{al-} is still on the front.
\end{cousinweb}

\begin{grammarlens}[title={What you've built}]
One more word for reading notices, forms and messages.
\end{grammarlens}

\subsection*{Your turn}
\begin{itemize}
\raggedright
  \item \emph{Say it:} \emph{el alquiler}
  \item \emph{Say it:} \emph{el alquiler}, once more
  \item \emph{Say it:} a sentence of your own with \emph{el alquiler}
  \item \emph{Recall:} say the Spanish for bored, then the Spanish for the saving, then say \emph{el alquiler} again
\end{itemize}

\begin{tcolorbox}[breakable,colback=teal!4,colframe=teal!35!black,title={Before you move on}]
What does el alquiler mean? (\textquotedblleft{}The rent; renting\textquotedblright{}.) Say it once more.
\end{tcolorbox}
\section[la multa]{la multa — the fine}

\label{lesson:ES-C477-multa}



\begin{quote}
Before the new one: say the Spanish for the saving, then the Spanish for the rent.
\end{quote}

\subsection*{You'll want to know: la multa}
\textbf{la multa} — \textquotedblleft{}the fine\textquotedblright{}.

\begin{quote}
Me pusieron una multa por aparcar mal.
\end{quote}

\begin{quote}
\emph{They gave me a fine for parking badly.}
\end{quote}

\emph{Poner una multa} — to give a fine. \emph{Aparcar} is the word in Spain; in Latin America people say \emph{estacionar}.

\begin{cousinweb}
\textbf{multa} $\leftarrow$ Latin \emph{multa}, \textquotedblleft{}a fine, a penalty\textquotedblright{} — a word of Roman law, taken from Latin by scholars. The Romans themselves called it a Sabine or an Oscan word, from two of their Italian neighbours (Lewis \& Short).
\end{cousinweb}

\begin{grammarlens}[title={What you've built}]
One more word for reading notices, forms and messages.
\end{grammarlens}

\subsection*{Your turn}
\begin{itemize}
\raggedright
  \item \emph{Say it:} \emph{la multa}
  \item \emph{Say it:} \emph{la multa}, once more
  \item \emph{Say it:} a sentence of your own with \emph{la multa}
  \item \emph{Recall:} say the Spanish for the saving, then the Spanish for the rent, then say \emph{la multa} again
\end{itemize}

\begin{tcolorbox}[breakable,colback=teal!4,colframe=teal!35!black,title={Before you move on}]
What does la multa mean? (\textquotedblleft{}The fine\textquotedblright{}.) Say it once more.
\end{tcolorbox}
\section[ganar]{ganar — to earn; to win}

\label{lesson:ES-C477-ganar}



\begin{quote}
Before the new one: say the Spanish for the rent, then the Spanish for the fine.
\end{quote}

\subsection*{You'll want to know: ganar}
\textbf{ganar} — \textquotedblleft{}to earn; to win\textquotedblright{}.

\begin{quote}
Gana mil euros al mes.
\end{quote}

\begin{quote}
\emph{She earns a thousand euros a month.}
\end{quote}

The same verb wins a game: \emph{ganar un partido}.

\begin{cousinweb}
\textbf{ganar} is probably Germanic. The usual account traces it to Gothic \emph{*ganan}, \textquotedblleft{}to covet\textquotedblright{} (EtymDB 2.1 records this line), perhaps crossed with the Germanic verb for grazing and hunting that is behind French \emph{gagner}.
\end{cousinweb}

\begin{grammarlens}[title={What you've built}]
One more word for reading notices, forms and messages.
\end{grammarlens}

\subsection*{Your turn}
\begin{itemize}
\raggedright
  \item \emph{Say it:} \emph{ganar}
  \item \emph{Say it:} \emph{ganar}, once more
  \item \emph{Say it:} a sentence of your own with \emph{ganar}
  \item \emph{Recall:} say the Spanish for the rent, then the Spanish for the fine, then say \emph{ganar} again
\end{itemize}

\begin{tcolorbox}[breakable,colback=teal!4,colframe=teal!35!black,title={Before you move on}]
What does ganar mean? (\textquotedblleft{}To earn; to win\textquotedblright{}.) Say it once more.
\end{tcolorbox}
\section[trescientos]{trescientos — three hundred}

\label{lesson:ES-C477-trescientos}



\begin{quote}
Before the new one: say the Spanish for the fine, then the Spanish for to earn.
\end{quote}

\subsection*{You'll want to know: trescientos}
\textbf{trescientos} — \textquotedblleft{}three hundred\textquotedblright{}.

\begin{quote}
El billete cuesta trescientos euros.
\end{quote}

\begin{quote}
\emph{The ticket costs three hundred euros.}
\end{quote}

It agrees with the noun: \emph{trescientas personas}.

\begin{cousinweb}
\textbf{trescientos} $\leftarrow$ Old Spanish \emph{trezientos} (EtymDB 2.1), from Latin \emph{trecenti}, \textquotedblleft{}three hundred\textquotedblright{} — inherited: the stressed \emph{e} opened into \emph{ie}. The modern spelling remade the front to match \emph{tres}.
\end{cousinweb}

\begin{grammarlens}[title={What you've built}]
One more word for reading notices, forms and messages.
\end{grammarlens}

\subsection*{Your turn}
\begin{itemize}
\raggedright
  \item \emph{Say it:} \emph{trescientos}
  \item \emph{Say it:} \emph{trescientos}, once more
  \item \emph{Say it:} a sentence of your own with \emph{trescientos}
  \item \emph{Recall:} say the Spanish for the fine, then the Spanish for to earn, then say \emph{trescientos} again
\end{itemize}

\begin{tcolorbox}[breakable,colback=teal!4,colframe=teal!35!black,title={Before you move on}]
What does trescientos mean? (\textquotedblleft{}Three hundred\textquotedblright{}.) Say it once more.
\end{tcolorbox}
